\section{Discussion}
\label{sec:discussion}

\subsection{What this changes for people building HPC coding agents}

Three recommendations follow directly from the measurements, and none of them
require adopting our reward.

\emph{Do not use output equivalence as an acceptance gate for parallel code.}
It is necessary and it is not sufficient, and the gap is not small: in our
accepted set it is the difference between \SpeedMin$\times$ and
\SpeedMax$\times$. If an agent's loop terminates when the output matches, it
will terminate on programs that are slower than the code they replaced.

\emph{Report the target machine alongside any performance reward.} Correctness
is portable; efficiency is not. A model post-trained against measured speedup
on one architecture has been aligned to that architecture. We measure on a
single machine and do not claim cross-architecture transfer; characterizing how
far the efficiency ranking of these programs moves across machines is the
natural next experiment, and one the released harness is built to run.

\emph{Report reward resolution.} If the noise floor of the timing harness
exceeds the differences the reward is meant to express, the reward is noise.
This is cheap to measure and, as far as we can tell, not currently reported by
anyone---ourselves included, until we looked.

\subsection{Why we do not claim a trained policy}

The natural completion of this work is to post-train under both rewards and
show that the correctness-only arm drifts toward correct-but-serial code while
the gated arm does not. We deliberately do not report such a run.

The reason is that we would rather state a falsifiable prediction than a
rushed training curve. The prediction is specific: under GRPO with
$\Rcorr$ on a task distribution with $p_t$ near saturation, mean sampled
efficiency should be flat or declining while correctness stays constant,
because Proposition~\ref{prop:collapse} says most groups contribute no
gradient and Proposition~\ref{prop:degeneracy} says the surviving gradient
does not point toward parallelism. Under $\Rgate$ on the same distribution,
correctness should be preserved and mean sampled efficiency should rise. The
harness we release computes both rewards over sampled completions and is
wired for this experiment; we state the prediction so that it can be checked
against us.

\subsection{Threats to validity}

\paragraph*{The proposition may look trivial}
It is elementary once stated, and we say so in Section~\ref{sec:degeneracy}.
Our defence is that it is not stated anywhere in the RLVR-for-HPC literature
we surveyed, that correctness-only verifiable rewards are being deployed for
this exact task, and that the corollary about signal collapse
(Proposition~\ref{prop:collapse}) and the measured cost
(Section~\ref{sec:results-bakeoff}) are not elementary.

\paragraph*{Prior art on the phenomenon}
PARACODEX~\cite{paracodex} documented correctness-passing code that executes
serially, in a GPU-offload setting, in January 2026. We do not claim priority
on the observation. We claim the reframing---from an agent and harness defect
to be patched, to the location of the reward function's optimum---and the CPU
thread-scaling instantiation.

\paragraph*{Benchmark difficulty}
A reviewer may object that correctness looks saturated only because
\ParEval's OpenMP tasks are idiomatic. We agree, and it is the argument: on
tasks where correctness saturates, a correctness-only reward has nothing left
to teach, while the efficiency axis remains wide open. The objection would
bite if the efficiency spread narrowed on easy tasks. It does not---the
\SpeedMin$\times$--\SpeedMax$\times$ range is measured on exactly these
saturated tasks.

\paragraph*{Best-of-$N$ is not GRPO}
Our selection analysis captures the selection component of group-relative
optimization and not its distribution-sharpening
effect~\cite{bonmaxk,rewardunlikely}. We label it as a proxy wherever it
appears and read it as a lower bound.

\paragraph*{Scope}
Shared-memory OpenMP, \NTasks\ \ParEval\ tasks, a single measurement machine, one
serial reference per task. We do not generalize to MPI, CUDA, or Kokkos,
although the argument in Section~\ref{sec:degeneracy} does not depend on the
programming model and the harness supports them. The hazard list underlying
\PG\ is not exhaustive, and we expect it to be extended---that is the point of
releasing it.

\paragraph*{The reference defines the task}
Each task has a single serial reference. Being sequential it is immune to the
bugs under study, but it also fixes what the task means; the scan-family
ambiguity in Section~\ref{sec:results-saturation} is a concrete instance of
that fixing being contestable.
